% Part: propositional-logic

\documentclass[../../include/open-logic-part]{subfiles}

\begin{document}

\olpart{pl}{Propositielogica}

\begin{editorial}
  Dit deel behandelt de klassieke propositielogica. Het eerste hoofdstuk
  is nog vrij elementair en bevat hoofdzakelijk een opsomming van
  definities en resultaten; veel bewijzen worden niet uitgewerkt maar
  als opgaven aan de lezer overgelaten. De tekst over bewijssystemen en
  de volledigheidsstelling is overgenomen uit het deel over
  eerste-ordelogica, waarbij de tag ``FOL'' op onwaar is gezet. Hierdoor
  vervalt alles wat betrekking heeft op predicaten, termen en kwantoren,
  en worden passages over structuren~$\Struct{M}$ vervangen door
  passages over valuaties~$\pAssign{v}$.

  Het is de bedoeling dit deel verder uit te werken en er meer
  onderwerpen en resultaten aan toe te voegen, zoals
  functionele volledigheid.
\end{editorial}

\olimport[syntax-and-semantics]{syntax-and-semantics}

\tagfalse{FOL}

\olimport[../first-order-logic/proof-systems]{proof-systems}

\iftag{prfSC}{%
  \olimport[../first-order-logic/sequent-calculus]{sequent-calculus}
}{}

\iftag{prfND}{%
  \olimport[../first-order-logic/natural-deduction]{natural-deduction}
}{}

\iftag{prfTab}{%
  \olimport[../first-order-logic/tableaux]{tableaux}
}{}

\iftag{prfAX}{%
  \olimport[../first-order-logic/axiomatic-deduction]{axiomatic-deduction}
}{}

\olimport[../first-order-logic/completeness]{completeness}

\tagtrue{FOL}

\OLEndPartHook
\end{document}
